% Pregunta 876 (2003-julio-OEP2003, n.º 4): convergencia de Cournot; al crecer N el equilibrio baja por la demanda hacia p = c y la pérdida de eficiencia se reduce hasta desaparecer. Autor: Víctor Gutiérrez Marcos
\documentclass[tikz,border=6pt]{standalone}
\input{tikz-tcee}
\begin{document}
\begin{tikzpicture}[tcee,x=1cm,y=1cm]
  % Demanda p = 4,6 - (2/3)Q ; c = 1 ; (a-c)/b = 5,4
  % N=1: Q = 2,7, p = 2,8 ; N=2: Q = 3,6, p = 2,2 ; N->inf: Q = 5,4, p = 1
  \fill[area] (2.7,2.8) -- (2.7,1) -- (3.6,1) -- (3.6,2.2) -- cycle;
  \fill[area2] (3.6,2.2) -- (3.6,1) -- (5.4,1) -- cycle;
  \ejes{$Q$}{$p$}{6.6}{5}
  \draw[curva1] (0,4.6) -- (6,0.6) node[etiqueta,below right] {$D$};
  \draw[curva2] (0,1) -- (6.2,1) node[etiqueta,right] {$CMg$};
  \draw[guia] (2.7,0) node[below] {$Q_1$} -- (2.7,2.8) -- (0,2.8) node[left] {$p_1$};
  \draw[guia] (3.6,0) node[below] {$Q_2$} -- (3.6,2.2) -- (0,2.2) node[left] {$p_2$};
  \draw[guia] (5.4,0) node[below] {$Q_c$} -- (5.4,1);
  \node[left] at (0,1) {$c$};
  \node[punto] at (2.7,2.8) {};
  \node[punto] at (3.6,2.2) {};
  \node[punto,tceerojo] at (5.4,1) {};
  \node[etiqueta,above right] at (2.7,2.8) {$N=1$};
  \node[etiqueta,above right] at (3.6,2.2) {$N=2$};
  \node[etiqueta,above right=8pt and 1pt,tceerojo] at (5.4,1) {$N\to\infty$};
  % Leyenda
  \fill[area2] (0.2,-1.0) rectangle (0.5,-0.75); \node[anchor=base west] at (0.55,-0.97) {pérdida de eficiencia con $N=2$};
  \fill[area] (0.2,-1.5) rectangle (0.5,-1.25); \fill[area2] (0.5,-1.5) rectangle (0.8,-1.25); \node[anchor=base west] at (0.85,-1.47) {pérdida con $N=1$ (las dos zonas)};
\end{tikzpicture}
\end{document}
